\scanpage{19}
\mspage{10}
% archive sheet number: 19
% lines 1--5
généraux que l'atlas canonique associé à un magasin -- et de revoir les prop. essentielles prop 1 à prop 5 (plus d'autres à venir peut-être) dans cet angle, \unsure{logiquement} \unsure{effleuré|effleurant|effarant} à peine.

% line 6
\uline{27 juin} (ou plutôt, 28 juin ! \unsure{l'on ne sait plus|on ne sait plus|(on ne sait plus}\ldots)

% lines 7--11
\uline{Définition}. Soit $\unsure{\mathscr{F}|\mathcal{F}},\leq$ un ensemble ordonné. On appelle \uline{antifiltre} de $\mathscr{F}$ une partie $\mathscr{F}'$ de $\mathscr{F}$ satisfaisant les conditions suivantes

a) $\mathscr{F}'$ est fermé pour $\leq$, i.e. $F\in\mathscr{F}'$, $G\leq F\Rightarrow G\in\mathscr{F}'$.

b) Si $F,G\in\mathscr{F}'$, et si $F\vee G$ existe dans $\mathscr{F}$, alors $F\vee G\in\mathscr{F}'$.

% lines 12--13, cancelled
\deleted{On dit que l'antifiltre est \uline{\unsure{couvrant}} si}
\[
\deleted{\unsure{\bigcup_{F\in\mathscr{F}}F=\Omega|\bigcup_{F\in\mathscr{F}}F=\mathscr{F}}}
\]

% lines 14--20
\uline{Définition}. Soit $M$ un magasin. Un \uline{atlas \unsure{attaché|associé} à $M$}, \inserted{ou $M$-atlas,} est par définition un antifiltre $\mathscr{F}$ (dans l'``ens'' $\mathscr{F}_M$ des \inserted{$M$-}figures, \unsure{inconditionnelles}, \deleted{tel que la figure vide et} contenant les \deleted{figure vide} figures élémentaires $\deleted{\unsure{X}}\widetilde{X}$ $(X\in M)$ et la figure vide \deleted{Cette dernière condition résulte des autres si $M\neq\varnothing$}).

% margin note written sideways
\margin{\uline{\unsure{Mag atl}} \inserted{sur $\mathscr{F}\subset\mathscr{F}_M$}\par
a) $F\in\mathscr{F}$, $G\leq F\Rightarrow G\in\mathscr{F}$.\par
b) $F,G\in\mathscr{F}$, $F\unsure{\gtrless|\bowtie}G\Rightarrow F\cup G\in\mathscr{F}$.\par
c) $\forall X\in M$, $\widetilde{X}\in\mathscr{F}$, et\par
d) $\mathscr{F}\neq\varnothing$ (i.e. $\varnothing\in\mathscr{F}$).}
% margin note written sideways
\margin{\unsure{revient à dire que} $\mathscr{F}$ (en tant que $\subset\unsure{\mathscr{P}|\mathcal{P}}(M)$) \uline{couvre} $M$ $\Longleftrightarrow$ c)}

% lines 21--26
\uline{NB} Cette dernière condition résulte des autres si $M\neq\varnothing$. On voit que les conditions suivantes sont équivalentes
\[
\begin{aligned}
\text{(i)}\quad&M=\varnothing\\
\text{(ii)}\quad&\mathscr{F}=\{\varnothing\}\\
\text{(iii)}\quad&\operatorname{Card}\mathscr{F}=1
\end{aligned}
\]
% lines 27--30
On dit \deleted{que} alors que \uline{le magasin} \deleted{$M$ est} \inserted{$M$ est} \uline{vide}, ou que \uline{l'atlas} \deleted{$\mathscr{F}$} \inserted{$\mathscr{F}$ est} \uline{vide} (un abus de langage pour $\mathscr{F}$, car $\mathscr{F}$ en tant qu'ens n'est pas vide \ldots.
